\scanpage{39}
\uline{Cor} Moyennant (R), $\mathcal X$ satisfait la condition
des chaines, et $\delta$ est l'unique fonction
$\mathcal X\mapsto\mathbb N$ telle que

a) $\delta(y)=\delta(x)+1$ si $y$ est un successeur de $x$

b) $\delta(x)=0$ si $x$ élément minimal de $\mathcal X$.

% Dark vertical emphasis mark at the left of the proposition.
\uline{Prop} \inserted{(Moyennant R)} Soit $x<y$, alors $|\mathcal X_{x<\cdot<y}|$ est une
sphère topologique de dim. $\delta(y)-\delta(x)-2$
(on convient de considérer l'espace vide comme
sphère top. de dim. $-1$)

Ceci résulte du

% Dark vertical emphasis mark at the left of the lemma.
\uline{Lemme} Soit $\mathcal Y$ un ens ordonné fini tel que
$|\mathcal Y|$ soit une $d$-variété topologique. Alors

a) $\forall x\in\mathcal Y$, $|\mathcal Y_{>x}|$ est une sphère top. de
dim $d-\delta(x)-2$

b) $\forall x\in\mathcal Y$, $|\mathcal Y_{<x}|$ est une sphère top.
de dim $\delta(x)-1$ (i.e. $\mathcal Y$ satisfait la
condition (R))

c) $\forall z,x\in\mathcal Y$, $z<x$, $|\mathcal Y_{z<\,\,<x}|$
\deleted{\unsure{satisfait|satisfai}} est une sphère topologique de dim
$\delta(x)-\delta(z)-2$.
% archive sheet number: 38
